\begin{table}[ht]
\centering
\small
\caption{Exact operating characteristic of the overshoot criterion of Equation~(\ref{eq:overshoot-threshold}), at $m' = 200$ exploration shots. Nothing here is simulated: the probe and its reference each have an enumerable transformed-binomial law, so the probabilities are computed by convolution over all $(m'+1)^2$ outcome pairs. $r = N\phi/(\pi/2)$ is the probe's depth relative to the first aliasing boundary and $\rho = N/N_{\mathrm{acc}}$ is how much deeper the probe is than the accepted estimate it is compared against, so small $\rho$ is a precise reference and large $\rho$ a noisy one. \emph{Detected from} is the smallest overshoot the rule catches with probability $0.9$, and the last column converts that blind spot into the relative bias $2(1-1/r)$ that an undetected overshoot of exactly that size would leave in the estimate. Note that the achieved false-alarm rate matches the nominal $\alpha$ only for an exact reference: replacing $\phi$ by a noisy $\hat\phi_{\mathrm{acc}}$ inflates it substantially, which is why $\alpha$ is tuned rather than set to a conventional confidence level.}
\label{tab:overshoot-operating}
\setlength{\tabcolsep}{5pt}
\begin{tabular}{lccc}
\toprule
Reference
& \makecell{False alarm \\ at $r=0.9$}
& \makecell{Detected from \\ $r \geq$}
& \makecell{Largest undetected \\ bias $|\phi-\bar\phi|/\phi$} \\
\midrule
\multicolumn{4}{l}{\itshape nominal level $\alpha = 0.5$} \\
exact ($\hat\phi_{\mathrm{acc}} = \phi$)
& 54.3\%
& 1.005
& 1.0\% \\
$\rho = 1.5$
& 48.5\%
& 1.034
& 6.6\% \\
$\rho = 3$
& 46.4\%
& 1.051
& 9.6\% \\
$\rho = 6$
& 47.3\%
& 1.095
& 17.4\% \\
\midrule
\multicolumn{4}{l}{\itshape nominal level $\alpha = 0.05$} \\
exact ($\hat\phi_{\mathrm{acc}} = \phi$)
& 2.6\%
& 1.039
& 7.5\% \\
$\rho = 1.5$
& 16.6\%
& 1.052
& 9.9\% \\
$\rho = 3$
& 28.0\%
& 1.070
& 13.0\% \\
$\rho = 6$
& 36.0\%
& 1.114
& 20.5\% \\
\bottomrule
\end{tabular}
\end{table}
